\textbf{Phần 1. Trắc nghiệm lựa chọn}
\setcounterpageref{ex}{0}

\begin{ex} %[2D10V53]
Điền dấu thích hợp vào chỗ chấm: $645 \dots 654$.
\choice
{$>$}
{\True $<$}
{$=$}
{Không so sánh được}
\loigiai{
Hàng trăm bằng nhau, hàng chục $4 < 5$ nên $645 < 654$.
\textbf{Đáp án của câu hỏi là B}
}
\end{ex}

\begin{ex} %[2D10V53]
Trong các số $789, 798, 879, 897$, số lớn nhất là
\choice
{\True $897$}
{$798$}
{$879$}
{$789$}
\loigiai{
So sánh hàng trăm, 8 lớn hơn 7. So sánh tiếp $897 > 879$.
\textbf{Đáp án của câu hỏi là A}
}
\end{ex}

\begin{ex} %[2D10V53]
Số nào dưới đây lớn hơn $456$ và bé hơn $462$?
\choice
{$452$}
{$465$}
{\True $459$}
{$470$}
\loigiai{
$456 < 459 < 462$.
\textbf{Đáp án của câu hỏi là C}
}
\end{ex}

\begin{ex} %[2D10V53]
Cho các số: $341, 314, 431, 413$. Thứ tự các số từ bé đến lớn là
\choice
{$431, 413, 341, 314$}
{$341, 314, 413, 431$}
{$314, 413, 341, 431$}
{\True $314, 341, 413, 431$}
\loigiai{
$314 < 341 < 413 < 431$.
\textbf{Đáp án của câu hỏi là D}
}
\end{ex}

\begin{ex} %[2D10V53]
Điền chữ số thích hợp vào dấu $*$ để: $5*8 > 588$.
\choice
{$8$}
{\True $9$}
{$7$}
{$0$}
\loigiai{
Để $5*8 > 588$ thì chữ số hàng chục $*$ phải lớn hơn 8, vậy $* = 9$.
\textbf{Đáp án của câu hỏi là B}
}
\end{ex}

\begin{ex} %[2D10V53]
Điền dấu thích hợp: $400 + 50 + 8 \dots 400 + 60$.
\choice
{\True $<$}
{$>$}
{$=$}
{Không so sánh được}
\loigiai{
$400 + 50 + 8 = 458$; $400 + 60 = 460$. Vì $458 < 460$ nên điền dấu $<$. 
\textbf{Đáp án của câu hỏi là A}
}
\end{ex}

\begin{ex} %[2D10V53]
Có bao nhiêu số tự nhiên $x$ có ba chữ số thỏa mãn: $198 < x < 203$?
\choice
{$3$ số}
{$5$ số}
{\True $4$ số}
{$2$ số}
\loigiai{
Các số là: 199, 200, 201, 202. Có tất cả 4 số.
\textbf{Đáp án của câu hỏi là C}
}
\end{ex}

\begin{ex} %[2D10V53]
Lớp 2A quyên góp được $325$ quyển vở, lớp 2B được $340$ quyển vở, lớp 2C được $318$ quyển vở. Lớp quyên góp được nhiều vở nhất là
\choice
{Lớp 2A}
{Lớp 2C}
{Cả ba lớp bằng nhau}
{\True Lớp 2B}
\loigiai{
$340 > 325 > 318$ nên lớp 2B quyên góp nhiều nhất.
\textbf{Đáp án của câu hỏi là D}
}
\end{ex}

\begin{ex} %[2D10V53]
Tìm số có ba chữ số bé nhất trong các số sau:
\choice
{$205$}
{\True $189$}
{$250$}
{$199$}
\loigiai{
So sánh hàng trăm 1 và 2, giữa 189 và 199 thì hàng chục $8 < 9$ nên 189 bé nhất.
\textbf{Đáp án của câu hỏi là B}
}
\end{ex}

\begin{ex} %[2D10V53]
Khoảng cách từ Hà Nội đến Vinh khoảng $300\text{ km}$, đến Thanh Hóa khoảng $160\text{ km}$, đến Huế khoảng $660\text{ km}$. Nơi nào xa Hà Nội nhất?
\choice
{\True Huế}
{Vinh}
{Thanh Hóa}
{Ba nơi xa như nhau}
\loigiai{
$660 > 300 > 160$ nên Huế xa Hà Nội nhất.
\textbf{Đáp án của câu hỏi là A}
}
\end{ex}

\begin{ex} %[2D10V53]
Cho số $A = 700 + 30 + 5$ và số $B = 739$. Khẳng định nào sau đây ĐÚNG?
\choice
{$A > B$}
{$A = B$}
{\True $A < B$}
{$A - B = 4$}
\loigiai{
$A = 735 < 739 = B$.
\textbf{Đáp án của câu hỏi là C}
}
\end{ex}

\begin{ex} %[2D10V53]
Dãy số nào sau đây gồm toàn các số lẻ và được xếp theo thứ tự giảm dần?
\choice
{$888, 777, 666$}
{$135, 357, 579$}
{$999, 998, 997$}
{\True $987, 875, 763$}
\loigiai{
Dãy $987, 875, 763$ gồm toàn số lẻ và giảm dần.
\textbf{Đáp án của câu hỏi là D}
}
\end{ex}

\textbf{Phần 2. Trắc nghiệm đúng sai}
\setcounterpageref{ex}{0}

\begin{ex} %[2D10V53]
Xét các mệnh đề so sánh số có ba chữ số:
\choiceTF
{\True Số nào có chữ số hàng trăm lớn hơn thì lớn hơn.}
{\True 999 là số có ba chữ số lớn nhất.}
{Số 542 bé hơn số 524.}
{\True Có duy nhất một số tự nhiên nằm giữa 699 và 701.}
\loigiai{
\begin{itemize}
    \item Ý a) Đúng: Khi so sánh 2 số có 3 chữ số, số có chữ số hàng trăm lớn hơn thì lớn hơn.
    \item Ý b) Đúng: Đúng, số liền sau là 1 000 có 4 chữ số.
    \item Ý c) Sai: Hàng chục $4 > 2$ nên $542 > 524$.
    \item Ý d) Đúng: Đó là số 700.
\end{itemize}
\begin{itemize}
    \item a) \textbf{Đúng}.
    \item b) \textbf{Đúng}.
    \item c) \textbf{Sai}.
    \item d) \textbf{Đúng}.
\end{itemize}
}
\end{ex}

\begin{ex} %[2D10V53]
Cho bốn số 425, 452, 524, 542:
\choiceTF
{\True Số bé nhất trong bốn số là 425.}
{\True Số lớn nhất trong bốn số là 542.}
{452 lớn hơn 524.}
{\True Hiệu giữa số lớn nhất và số bé nhất là 117.}
\loigiai{
\begin{itemize}
    \item Ý a) Đúng: Đúng.
    \item Ý b) Đúng: Đúng.
    \item Ý c) Sai: Hàng trăm $4 < 5$ nên $452 < 524$.
    \item Ý d) Đúng: 542 - 425 = 117. Đúng.
\end{itemize}
\begin{itemize}
    \item a) \textbf{Đúng}.
    \item b) \textbf{Đúng}.
    \item c) \textbf{Sai}.
    \item d) \textbf{Đúng}.
\end{itemize}
}
\end{ex}

\textbf{Phần 3. Trả lời ngắn}
\setcounterpageref{ex}{0}

\begin{ex} %[2D10V53]
Cho bốn chữ số 3, 0, 8, 5. Viết số chẵn lớn nhất có ba chữ số khác nhau.
\shortans{850}
\loigiai{
\textbf{Đáp số:} $850$\\
Hàng trăm 8, chục 5, đơn vị chẵn 0: số 850.
}
\end{ex}

\begin{ex} %[2D10V53]
Có bao nhiêu số có ba chữ số lớn hơn 895 và bé hơn 905?
\shortans{9}
\loigiai{
\textbf{Đáp số:} $9$\\
Từ 896 đến 904 có: $904 - 896 + 1 = 9$ số.
}
\end{ex}

\begin{ex} %[2D10V53]
Tìm hiệu giữa số lớn nhất có ba chữ số khác nhau và số bé nhất có ba chữ số khác nhau.
\shortans{885}
\loigiai{
\textbf{Đáp số:} $885$\\
Số lớn nhất: 987. Số bé nhất: 102. Hiệu: $987 - 102 = 885$.
}
\end{ex}

\begin{ex} %[2D10V53]
Tìm chữ số thích hợp thay cho $*$ để $478 < 4*8 < 498$.
\shortans{8}
\loigiai{
\textbf{Đáp số:} $8$\\
Hàng chục $*$ phải lớn hơn 7 và bé hơn 9 nên $* = 8$.
}
\end{ex}

\begin{ex} %[2D10V53]
Sắp xếp các số 502, 205, 520, 250 theo thứ tự tăng dần. Số thứ hai trong dãy là số nào?
\shortans{250}
\loigiai{
\textbf{Đáp số:} $250$\\
Thứ tự: 205, 250, 502, 520. Số thứ hai là 250.
}
\end{ex}

\begin{ex} %[2D10V53]
Tìm số tự nhiên liền trước của số có ba chữ số bé nhất.
\shortans{99}
\loigiai{
\textbf{Đáp số:} $99$\\
Số bé nhất có 3 chữ số là 100. Liền trước là 99.
}
\end{ex}
